\IfFileExists{stacks-project.cls}{%
\documentclass{stacks-project}
}{%
\documentclass{amsart}
}
\usepackage{amssymb}

% For dealing with references we use the comment environment
\usepackage{verbatim}
\usepackage{fancyvrb}
\newenvironment{reference}{\comment}{\endcomment}
%\newenvironment{reference}{}{}
\newenvironment{slogan}{\comment}{\endcomment}
\newenvironment{history}{\comment}{\endcomment}

% For source figures retained as assets
\usepackage{graphicx}

% For commutative diagrams we use Xy-pic
\usepackage[all]{xy}

% We use 2cell for 2-commutative diagrams.
\xyoption{2cell}
\UseAllTwocells

% We use multicol for the list of chapters between chapters
\usepackage{multicol}

% This is generally recommended for better output
\usepackage{lmodern}
\usepackage[T1]{fontenc}
\usepackage{textalpha}

% For cross-file-references
% We prefer the snapshot tmp/xref/#2.aux,
% which latexmkrc replaces atomically, so
% that documents compiled simultaneously
% never read a half-written .aux file.
% The last argument keeps links to #2.pdf.
\usepackage{xr-hyper}
\newcommand{\maybeexternaldocument}[2]{%
  \IfFileExists{#2.aux}{%
    \externaldocument[#1]{#2}%
  }{%
    \PackageWarning{stack}{Missing #2.aux; external references with prefix #1 unavailable}%
  }%
}
\let\externaldocumentorig\externaldocument
\renewcommand{\externaldocument}[2][]{%
  \IfFileExists{tmp/xref/#2.aux}{%
    \externaldocumentorig[#1]{tmp/xref/#2}[#2.pdf]%
  }{%
  \IfFileExists{#2.aux}{%
    \externaldocumentorig[#1]{#2}%
  }{%
    \PackageWarning{stack}{Missing #2.aux; external references with prefix #1 unavailable}%
  }%
  }%
}

% Package for hypertext links:
\usepackage{hyperref}

% For any local file, say "hello.tex" you want to link to please
% use \externaldocument[hello-]{hello}
\externaldocument[introduction-]{introduction}
\externaldocument[conventions-]{conventions}
\externaldocument[sets-]{sets}
\externaldocument[categories-]{categories}
\externaldocument[topology-]{topology}
\externaldocument[sheaves-]{sheaves}
\externaldocument[sites-]{sites}
\externaldocument[stacks-]{stacks}
\externaldocument[fields-]{fields}
\externaldocument[algebra-]{algebra}
\externaldocument[brauer-]{brauer}
\externaldocument[homology-]{homology}
\externaldocument[derived-]{derived}
\externaldocument[simplicial-]{simplicial}
\externaldocument[more-algebra-]{more-algebra}
\externaldocument[smoothing-]{smoothing}
\externaldocument[modules-]{modules}
\externaldocument[sites-modules-]{sites-modules}
\externaldocument[injectives-]{injectives}
\externaldocument[cohomology-]{cohomology}
\externaldocument[sites-cohomology-]{sites-cohomology}
\externaldocument[dga-]{dga}
\externaldocument[dpa-]{dpa}
\externaldocument[sdga-]{sdga}
\externaldocument[hypercovering-]{hypercovering}
\externaldocument[schemes-]{schemes}
\externaldocument[constructions-]{constructions}
\externaldocument[properties-]{properties}
\externaldocument[morphisms-]{morphisms}
\externaldocument[coherent-]{coherent}
\externaldocument[divisors-]{divisors}
\externaldocument[limits-]{limits}
\externaldocument[varieties-]{varieties}
\externaldocument[topologies-]{topologies}
\externaldocument[descent-]{descent}
\externaldocument[perfect-]{perfect}
\externaldocument[more-morphisms-]{more-morphisms}
\externaldocument[flat-]{flat}
\externaldocument[groupoids-]{groupoids}
\externaldocument[more-groupoids-]{more-groupoids}
\externaldocument[etale-]{etale}
\externaldocument[chow-]{chow}
\externaldocument[intersection-]{intersection}
\externaldocument[pic-]{pic}
\externaldocument[weil-]{weil}
\externaldocument[adequate-]{adequate}
\externaldocument[dualizing-]{dualizing}
\externaldocument[duality-]{duality}
\externaldocument[discriminant-]{discriminant}
\externaldocument[derham-]{derham}
\externaldocument[local-cohomology-]{local-cohomology}
\externaldocument[algebraization-]{algebraization}
\externaldocument[curves-]{curves}
\externaldocument[resolve-]{resolve}
\externaldocument[models-]{models}
\externaldocument[functors-]{functors}
\externaldocument[equiv-]{equiv}
\externaldocument[pione-]{pione}
\externaldocument[etale-cohomology-]{etale-cohomology}
\externaldocument[proetale-]{proetale}
\externaldocument[relative-cycles-]{relative-cycles}
\externaldocument[more-etale-]{more-etale}
\externaldocument[trace-]{trace}
\externaldocument[crystalline-]{crystalline}
\externaldocument[spaces-]{spaces}
\externaldocument[spaces-properties-]{spaces-properties}
\externaldocument[spaces-morphisms-]{spaces-morphisms}
\externaldocument[decent-spaces-]{decent-spaces}
\externaldocument[spaces-cohomology-]{spaces-cohomology}
\externaldocument[spaces-limits-]{spaces-limits}
\externaldocument[spaces-divisors-]{spaces-divisors}
\externaldocument[spaces-over-fields-]{spaces-over-fields}
\externaldocument[spaces-topologies-]{spaces-topologies}
\externaldocument[spaces-descent-]{spaces-descent}
\externaldocument[spaces-perfect-]{spaces-perfect}
\externaldocument[spaces-more-morphisms-]{spaces-more-morphisms}
\externaldocument[spaces-flat-]{spaces-flat}
\externaldocument[spaces-groupoids-]{spaces-groupoids}
\externaldocument[spaces-more-groupoids-]{spaces-more-groupoids}
\externaldocument[bootstrap-]{bootstrap}
\externaldocument[spaces-pushouts-]{spaces-pushouts}
\externaldocument[spaces-chow-]{spaces-chow}
\externaldocument[groupoids-quotients-]{groupoids-quotients}
\externaldocument[spaces-more-cohomology-]{spaces-more-cohomology}
\externaldocument[spaces-simplicial-]{spaces-simplicial}
\externaldocument[spaces-duality-]{spaces-duality}
\externaldocument[formal-spaces-]{formal-spaces}
\externaldocument[restricted-]{restricted}
\externaldocument[spaces-resolve-]{spaces-resolve}
\externaldocument[formal-defos-]{formal-defos}
\externaldocument[defos-]{defos}
\externaldocument[cotangent-]{cotangent}
\externaldocument[examples-defos-]{examples-defos}
\externaldocument[algebraic-]{algebraic}
\externaldocument[examples-stacks-]{examples-stacks}
\externaldocument[stacks-sheaves-]{stacks-sheaves}
\externaldocument[criteria-]{criteria}
\externaldocument[artin-]{artin}
\externaldocument[quot-]{quot}
\externaldocument[stacks-properties-]{stacks-properties}
\externaldocument[stacks-morphisms-]{stacks-morphisms}
\externaldocument[stacks-limits-]{stacks-limits}
\externaldocument[stacks-cohomology-]{stacks-cohomology}
\externaldocument[stacks-perfect-]{stacks-perfect}
\externaldocument[stacks-introduction-]{stacks-introduction}
\externaldocument[stacks-more-morphisms-]{stacks-more-morphisms}
\externaldocument[stacks-geometry-]{stacks-geometry}
\externaldocument[moduli-]{moduli}
\externaldocument[moduli-curves-]{moduli-curves}
\externaldocument[examples-]{examples}
\externaldocument[exercises-]{exercises}
\externaldocument[guide-]{guide}
\externaldocument[desirables-]{desirables}
\externaldocument[coding-]{coding}
\externaldocument[obsolete-]{obsolete}
\externaldocument[fdl-]{fdl}
\externaldocument[index-]{index}

% Heap Project documents
\externaldocument[k3-]{k3}
\externaldocument[torelli-]{torelli}
\externaldocument[linear-algebra-]{linear-algebra}
\externaldocument[groups-]{groups}
\externaldocument[complex-analysis-]{complex-analysis}
\externaldocument[functional-analysis-]{functional-analysis}
\externaldocument[categories-]{categories}
\externaldocument[commutative-algebra-]{commutative-algebra}
\externaldocument[ringed-spaces-]{ringed-spaces}
\externaldocument[homological-algebra-]{homological-algebra}
\externaldocument[lie-algebras-]{lie-algebras}
\externaldocument[representation-theory-]{representation-theory}
\externaldocument[smooth-manifolds-]{smooth-manifolds}
\externaldocument[differential-topology-]{differential-topology}
\externaldocument[algebraic-topology-]{algebraic-topology}
\externaldocument[differential-geometry-]{differential-geometry}
\externaldocument[algebraic-geometry-]{algebraic-geometry}
\externaldocument[symplectic-geometry-]{symplectic-geometry}
\externaldocument[complex-geometry-]{complex-geometry}
\externaldocument[hodge-theory-]{hodge-theory}
\externaldocument[%
variations-of-hodge-structures-]{%
variations-of-hodge-structures}
\externaldocument[%
homological-mirror-symmetry-]{%
homological-mirror-symmetry}
\maybeexternaldocument{deformations-}{deformations}
\externaldocument[vertex-algebras-]{vertex-algebras}
\externaldocument[springer-]{springer}
\externaldocument[infty-categories-]{infty-categories}
\externaldocument[seminars-differential-geometry-]{%
seminars-differential-geometry}
\externaldocument[seminars-complex-geometry-]{%
seminars-complex-geometry}
\externaldocument[seminars-algebraic-geometry-]{%
seminars-algebraic-geometry}
\externaldocument[seminars-others-]{%
seminars-others}
\externaldocument[%
geometric-prequantization-]{%
geometric-prequantization}
\externaldocument[surfaces-]{surfaces}
\externaldocument[birational-maps-commutative-algebra-]{birational-maps-commutative-algebra}
\externaldocument[cremona-transformations-]{cremona-transformations}
\externaldocument[derived-categories-of-sheaves-]{derived-categories-of-sheaves}
\externaldocument[geometric-invariant-theory-]{geometric-invariant-theory}
\externaldocument[geometric-stability-conditions-and-group-actions-]{geometric-stability-conditions-and-group-actions}
\externaldocument[moduli-spaces-of-sheaves-]{moduli-spaces-of-sheaves}
\externaldocument[bridgeland-stability-general-theory-]{bridgeland-stability-general-theory}
\externaldocument[hilbert-schemes-]{hilbert-schemes}
\externaldocument[noncommutative-abelian-surfaces-and-kummer-type-hyperkahler-varieties-]{noncommutative-abelian-surfaces-and-kummer-type-hyperkahler-varieties}
\externaldocument[mumford-tate-groups-in-hodge-theory-]{mumford-tate-groups-in-hodge-theory}
\externaldocument[hodge-birational-atoms-2026-]{hodge-birational-atoms-2026}
\externaldocument[xxii-egd-2026-]{%
xxii-egd-2026}
\externaldocument[pde-]{pde}
\externaldocument[probability-]{probability}

\externaldocument[linguistics-]{linguistics}
\externaldocument[genealogia-familiar-]{%
genealogia-familiar}

\externaldocument[%
16th-alga-meeting-2026-]{%
16th-alga-meeting-2026}

\externaldocument[%
iv-jornadas-lefschetz-2026-]{%
iv-jornadas-lefschetz-2026}

\externaldocument[reading-the-masters-]{%
reading-the-masters}

\externaldocument[real-analysis-]{%
real-analysis}

\externaldocument[optimal-transport-]{%
optimal-transport}

\externaldocument[history-]{history}

\externaldocument[seiberg-witten-]{seiberg-witten}

\externaldocument[history-philosophy-biology-]{history-philosophy-biology}

\externaldocument[chemistry-]{chemistry}

\externaldocument[molecular-cell-biology-]{molecular-cell-biology}

\externaldocument[microbiology-]{microbiology}

\externaldocument[protistology-]{protistology}

\externaldocument[mycology-]{mycology}

\externaldocument[botany-]{botany}

\externaldocument[zoology-]{zoology}

\externaldocument[animal-histology-]{animal-histology}

\externaldocument[animal-physiology-]{animal-physiology}

\externaldocument[systematics-biogeography-]{systematics-biogeography}

\externaldocument[ecology-natural-resources-]{ecology-natural-resources}

\externaldocument[earth-sciences-]{earth-sciences}

\externaldocument[biostatistics-]{biostatistics}

\externaldocument[mathematical-systems-biology-]{mathematical-systems-biology}

\externaldocument[cimpa-el-salvador-2024-]{cimpa-el-salvador-2024}

% Theorem environments.
%
\theoremstyle{plain}
\newtheorem{theorem}[subsection]{Theorem}
\newtheorem{proposition}[subsection]{Proposition}
\newtheorem{lemma}[subsection]{Lemma}
\newtheorem{corollary}[subsection]{Corollary}

\theoremstyle{definition}
\newtheorem{definition}[subsection]{Definition}
\newtheorem{example}[subsection]{Example}
\newtheorem{exercise}[subsection]{Exercise}
\newtheorem{situation}[subsection]{Situation}

\theoremstyle{remark}
\newtheorem{remark}[subsection]{Remark}
\newtheorem{remarks}[subsection]{Remarks}

\numberwithin{equation}{subsection}

% Macros
%
\def\lim{\mathop{\mathrm{lim}}\nolimits}
\def\colim{\mathop{\mathrm{colim}}\nolimits}
\def\Spec{\mathop{\mathrm{Spec}}}
\def\Hom{\mathop{\mathrm{Hom}}\nolimits}
\def\Ext{\mathop{\mathrm{Ext}}\nolimits}
\def\SheafHom{\mathop{\mathcal{H}\!\mathit{om}}\nolimits}
\def\SheafExt{\mathop{\mathcal{E}\!\mathit{xt}}\nolimits}
\def\Sch{\mathit{Sch}}
\def\Mor{\mathop{\mathrm{Mor}}\nolimits}
\def\Ob{\mathop{\mathrm{Ob}}\nolimits}
\def\Sh{\mathop{\mathit{Sh}}\nolimits}
\def\NL{\mathop{N\!L}\nolimits}
\def\CH{\mathop{\mathrm{CH}}\nolimits}
\def\proetale{{pro\text{-}\acute{e}tale}}
\def\etale{{\acute{e}tale}}
\def\QCoh{\mathit{QCoh}}
\def\Ker{\mathop{\mathrm{Ker}}}
\def\Im{\mathop{\mathrm{Im}}}
\def\Coker{\mathop{\mathrm{Coker}}}
\def\Coim{\mathop{\mathrm{Coim}}}

% Boxtimes
%
\DeclareMathSymbol{\boxtimes}{\mathbin}{AMSa}{"02}

%
% Macros for moduli stacks/spaces
%
\def\QCohstack{\mathcal{QC}\!\mathit{oh}}
\def\Cohstack{\mathcal{C}\!\mathit{oh}}
\def\Spacesstack{\mathcal{S}\!\mathit{paces}}
\def\Quotfunctor{\mathrm{Quot}}
\def\Hilbfunctor{\mathrm{Hilb}}
\def\Curvesstack{\mathcal{C}\!\mathit{urves}}
\def\Polarizedstack{\mathcal{P}\!\mathit{olarized}}
\def\Complexesstack{\mathcal{C}\!\mathit{omplexes}}
% \Pic is the operator that assigns to X its picard group, usage \Pic(X)
% \Picardstack_{X/B} denotes the Picard stack of X over B
% \Picardfunctor_{X/B} denotes the Picard functor of X over B
\def\Pic{\mathop{\mathrm{Pic}}\nolimits}
\def\Picardstack{\mathcal{P}\!\mathit{ic}}
\def\Picardfunctor{\mathrm{Pic}}
\def\Deformationcategory{\mathcal{D}\!\mathit{ef}}


\begin{document}

\title{Protistology}
\maketitle

\phantomsection
\label{section-phantom}

\tableofcontents

\section{What this chapter covers}
\label{section-overview}

Protists are the eukaryotes that are not
plants, animals or fungi: an artificial,
polyphyletic collection, united mainly by
lacking the tissues of those three groups
\cite[Chapter 25]{willey2008prescott}.
Traditionally the photosynthetic ones are
called algae and the heterotrophic ones
protozoa; the cyanobacteria, although
prokaryotes, are studied as algae too. This
chapter covers what algae are and how they
differ from plants, the origin of their
plastids by endosymbiosis, the organisms on
the boundaries between bacteria, algae,
plants and animals, the algal divisions,
their levels of organization and life cycles,
the protozoa and their classification, the
protozoa that parasitize humans, and the
course's field practical.

\textbf{Sources of these notes.} The course
was
\emph{Biolog\'ia de Protistas y Algas},
second semester of the Biology degree at the
Facultad de Ciencias, UNAM, in the spring of
2014 (UNAM semester 2014-2). What survives is
Daniel's study material: an outline of all
the algal divisions (pigments, reserves,
organization, reproduction, ecology), notes
on levels of organization, reproduction and
life cycles, questionnaires on algae and on
green algae, introductions to practicals on
diatoms, on cyanobacteria and on protozoa
associated with a toad, the team report of a
field practical at the Cantera Oriente, his
talk ``Protozoos par\'asitos del hombre''
(May 2014), an exam table on free-living
amoebae, the class tables of protozoan phyla
in Corliss's 1994 classification, and
readings such as Margulis and colleagues on
the origin of the nucleus and S\'anchez-Vega
and colleagues on intestinal parasites in
Mexico City \cite{sanchezvega2000frecuencia}.
The notes were checked and completed with
Lee's \emph{Phycology}
\cite{lee2008phycology}, Raven's
\emph{Biology of Plants}
\cite{evert2013raven}, Prescott's
microbiology \cite{willey2008prescott},
Tortora's \cite{tortora2013microbiology} and
Hickman's \emph{Animal Diversity}
\cite{hickman2002animal}.

\section{What algae are}
\label{section-algae-plants}

\begin{definition}
\label{definition-alga}
\emph{Algae} are thallophytes (organisms
without roots, stems and leaves) that have
chlorophyll a as their primary photosynthetic
pigment and lack a sterile covering of cells
around their reproductive cells \cite[Chapter
1]{lee2008phycology}. A \emph{thallus} is a
body that is not differentiated into such
organs.
\end{definition}

This is a definition by level of
organization, not by descent: it gathers the
cyanobacteria, which are prokaryotes, and
several unrelated lineages of eukaryotes
\cite[Chapter 1]{lee2008phycology}. ``Algae''
has no phylogenetic meaning \cite[Chapter
25]{willey2008prescott}.

Algae against land plants. Even large algae
such as kelps, with a holdfast, a stalk
(stipe) and blades, are thalli: their parts
are not roots, stems and leaves with vascular
tissue. Land plants (bryophytes and vascular
plants, the embryophytes) share a set of
features that no alga has \cite[Chapter
16]{evert2013raven}:
\begin{enumerate}
\item gametangia (antheridia and archegonia)
with a protective sterile jacket of cells;
\item retention of the zygote and of the
young multicellular embryo inside the
archegonium;
\item a multicellular diploid sporophyte;
\item multicellular sporangia with a sterile
jacket;
\item spores with walls of sporopollenin,
which resists decay and drying;
\item tissues produced by an apical meristem.
\end{enumerate}
Where algae and plants agree is in
photosynthesis: the green algae have the same
pigments (chlorophylls a and b), the same
reserve (starch inside the plastid) and
cellulose walls, and land plants evolved from
among them
(Section~\ref{section-selected-groups}).

\section{Plastids and endosymbiosis}
\label{section-endosymbiosis}

Plastids descend from cyanobacteria. In
\emph{primary endosymbiosis}, a phagotrophic
eukaryote took a cyanobacterium into a food
vesicle and, instead of digesting it, kept it
as an endosymbiont that shared its
photosynthate \cite[Part III and Chapter
3]{lee2008phycology}. The symbiont lost its
wall and most of its genes (many moved to the
host nucleus); the food vesicle membrane and
the cyanobacterial plasma membrane became the
two membranes of the chloroplast envelope.
Mereschkowsky proposed the idea in 1905. The
evidence: plastids and cyanobacteria are of
similar size, evolve oxygen, have 70S
ribosomes, circular DNA without histones and
chlorophyll a, and their rRNA sequences are
similar. Sequence data suggest that all
plastids go back to a single primary
endosymbiosis. Plastids with two envelope
membranes, from this primary event, are found
in the Glaucophyta, Rhodophyta and
Chlorophyta (and land plants), grouped today
as the Archaeplastida \cite[Chapter
25]{willey2008prescott}.

In \emph{secondary endosymbiosis}, a
eukaryote engulfed a eukaryotic alga and kept
it, so the plastid is wrapped in extra
membranes, called chloroplast endoplasmic
reticulum \cite[Parts IV and
V]{lee2008phycology}:
\begin{itemize}
\item one membrane of chloroplast ER (three
membranes in all): Euglenophyta, whose
plastid probably came from a green alga, and
Dinophyta and the apicomplexans, probably
from a red alga;
\item two membranes of chloroplast ER (four
in all): Chlorarachniophyta (from a green
alga) and Cryptophyta (from a red alga),
which still keep a \emph{nucleomorph}, the
shrunken nucleus of the engulfed alga,
between the second and third membranes; and
the Heterokontophyta and Haptophyta, in which
the nucleomorph has been lost. Two membranes
of chloroplast ER evolved at least three
times.
\end{itemize}

The course also read Margulis and colleagues'
symbiogenetic hypothesis that the nucleus
itself arose from a merger of an
archaebacterium with a spirochete-like
bacterium; unlike the origin of plastids and
mitochondria, this remains a hypothesis.

\section{Organisms on the boundaries}
\label{section-boundaries}

Between bacteria and plants.
\begin{itemize}
\item Cyanobacteria are prokaryotes with
plant-like oxygenic photosynthesis
\cite[Chapter 2]{lee2008phycology}. Their
thylakoids lie free in the cytoplasm and
carry chlorophyll a and phycobilisomes,
aggregates of phycobiliproteins. Many
filamentous species have heterocysts,
thick-walled cells without oxygen-producing
photosynthesis in which nitrogenase fixes
$\mathrm{N_2}$; akinetes, resting cells full
of reserves; and hormogonia, short motile
filaments. Their stromatolites are about 2700
million years old, and their photosynthesis
built up the oxygen of the atmosphere
\cite[Chapter 1]{lee2008phycology}.
\item Prochlorophytes (\emph{Prochloron}, a
symbiont of ascidians; \emph{Prochlorothrix};
\emph{Prochlorococcus}, the smallest known
photosynthetic organism) are cyanobacteria
with chlorophyll b. They were once put in a
division Prochlorophyta and seen as the
ancestors of green plastids, but rRNA trees
scatter them among the other cyanobacteria:
chlorophyll b arose several times
\cite[Chapter 2]{lee2008phycology}
\cite[Chapter 13]{evert2013raven}.
\item Glaucophytes (\emph{Cyanophora},
\emph{Glaucocystis}) are eukaryotes whose
plastids, called cyanelles, still look like
endosymbiotic cyanobacteria: they have
phycobiliproteins and, unlike chloroplasts, a
layer of peptidoglycan between their two
membranes \cite[Chapter 3]{lee2008phycology}
\cite[Chapter 25]{willey2008prescott}. They
are seen as a living picture of an early
stage in the evolution of the chloroplast.
\end{itemize}

Between algae and animals. \emph{Euglena} is
claimed by both botanists (Raven treats the
euglenoids as algae) and zoologists (Hickman
as protozoa, phylum Euglenozoa, next to the
trypanosomes) \cite[Chapter
15]{evert2013raven} \cite[Chapter
5]{hickman2002animal}: it photosynthesizes
but swims, and two thirds of the euglenoid
genera are colourless heterotrophs.

Between protists and animals.
Choanoflagellates (\emph{Monosiga},
\emph{Salpingoeca}, \emph{Codosiga}) are
solitary or colonial flagellates with a
single flagellum surrounded by a collar of
microvilli, with which they catch food
particles \cite[Chapter
25]{willey2008prescott}. They look like the
collar cells (choanocytes) of sponges, and
rRNA sequences support a common ancestor of
choanoflagellates and animals \cite[Chapter
6]{hickman2002animal}. With the fungi they
form the Opisthokonta, eukaryotes whose
flagellated cells have one posterior
flagellum.

\section{The algal divisions}
\label{section-divisions}

The table follows the course's divisions and
Lee's grouping by plastid membranes
\cite[Chapter 1 and Part
V]{lee2008phycology}. ``Chl'' stands for
chlorophyll and ``phycobilins'' for the
phycobiliproteins; the reserve is in the
cytoplasm unless stated.

\begin{center}
\begin{tabular}{p{0.21\linewidth}
p{0.23\linewidth}p{0.18\linewidth}
p{0.26\linewidth}}
Division & Pigments & Reserve &
Flagella and features \\
\hline
Cyanobacteria & chl a, phycobilins &
glycogen (cyanophycean starch),
cyanophycin & none; prokaryotic;
heterocysts \\
Glaucophyta & chl a, phycobilins &
starch & two, in flagellate forms;
cyanelles \\
Rhodophyta & chl a, phycobilins &
floridean starch & none; pit
connections; mostly marine \\
Chlorophyta & chl a and b & starch
inside the plastid & usually equal,
smooth; cellulose wall \\
Eugleno\-phyta & chl a and b &
paramylon & one emergent, with fine
hairs; pellicle, no wall; red eyespot \\
Dinophyta & chl a and c, peridinin &
starch & transverse and longitudinal;
thecal plates; red tides \\
Heterokonto\-phyta & chl a and c,
fucoxanthin & chryso\-laminarin,
laminarin & front tinsel and back
smooth; diatoms, kelps \\
Haptophyta & chl a and c,
fucoxanthin & chryso\-laminarin & two
smooth and a haptonema; coccoliths \\
Cryptophyta & chl a and c, phycobilins
& starch & two, with hairs; periplast,
nucleomorph \\
Chlorarachnio\-phyta & chl a and b &
in a vesicle on the pyrenoid &
amoeboid; zoospores with one;
nucleomorph \\
\end{tabular}
\end{center}

Notes on the table, from Lee's chapters and
Daniel's outline: cyanobacteria and red algae
both lack flagella and have
phycobiliproteins, and red algal plastids,
like those of green algae, have only the two
envelope membranes; red algae store floridean
starch outside the plastid; the heterokont
flagella are a long anterior flagellum with
tripartite tubular hairs and a short smooth
posterior one; the haptonema is a thread-like
appendage between the two flagella of
haptophytes, different in structure from a
flagellum \cite[Chapter
22]{lee2008phycology}. Heterokontophyta
include the golden algae, yellow-green algae
(Xanthophyceae), diatoms (Bacillariophyceae)
and brown algae (Phaeophyceae).

\section{Notes on selected groups}
\label{section-selected-groups}

Euglena, the green cell with a red eye
\cite[Chapter 15]{evert2013raven}.
\emph{Euglena} is a unicellular freshwater
flagellate without a cell wall. Under its
plasma membrane runs a pellicle of helical
protein strips, flexible enough to let the
cell change shape. A long flagellum emerges
from an anterior pocket (the reservoir), and
a second one stays inside it. Next to the
swelling at the base of the emergent
flagellum lies a red eyespot (stigma) of
carotenoid droplets, free in the cytosol;
together they form the light-sensing system
with which the cell swims towards light. A
contractile vacuole expels excess water
through the reservoir. The plastids have
chlorophylls a and b, like those of green
algae, from which they probably came; the
reserve is paramylon, a glucan stored in the
cytosol. Euglenoids divide lengthwise while
swimming; no sexual reproduction is known.
Green algae such as \emph{Chlamydomonas} also
have an orange-red eyespot, but inside the
chloroplast; in euglenoids and
dinoflagellates it lies outside \cite[Chapter
1]{lee2008phycology}.

Dinoflagellates and red tides.
Dinoflagellates have two flagella, one in a
transverse groove (the girdle or cingulum)
and one in a longitudinal groove (the
sulcus), and many are armoured with cellulose
plates in vesicles under the membrane (an
upper epitheca and a lower hypotheca). Their
blooms colour the water red. Some produce
potent toxins that accumulate in shellfish
and cause paralytic shellfish poisoning in
people who eat them \cite[Chapter
7]{lee2008phycology}; examples are
\emph{Alexandrium} \cite[Chapter
12]{tortora2013microbiology} and
\emph{Karenia brevis}, the cause of red tides
off Florida \cite[Chapter
15]{evert2013raven}. Other dinoflagellates
are the zooxanthellae, symbionts that feed
reef-building corals and giant clams.

Diatoms. Their wall, the frustule, is made of
silica and has two halves that fit like a
Petri dish; centric diatoms are radially
symmetric, pennate ones bilaterally. Their
plastids have chlorophylls a, c1 and c2 and
the golden-brown fucoxanthin \cite[Chapter
17]{lee2008phycology}. They are major members
of plankton and of the films on submerged
surfaces, and their fossil frustules form
deposits of diatomaceous earth.

Brown algae and kelps. All are multicellular,
almost all marine; kelps reach about 30 m,
with a holdfast, stipe, blades and often
gas-filled bladders, and their walls give the
alginates used as thickeners \cite[Chapter
15]{evert2013raven}. The reserve is
laminarin.

Red algae. They have no flagellated stages,
and their cells are joined by pit connections
\cite[Chapter 4]{lee2008phycology}. Most have
a life history with three phases: a haploid
gametophyte, a diploid carposporophyte that
grows on the female gametophyte and releases
carpospores, and a free-living diploid
tetrasporophyte whose meiosis gives
tetraspores \cite[Chapter
15]{evert2013raven}.

Green algae and the origin of land plants.
Green algae split into a chlorophyte clade
(most green algae) and a streptophyte clade,
which contains the charophycean algae and the
land plants \cite[Chapters 15 and
16]{evert2013raven}. Among charophytes,
\emph{Coleochaete} and the stoneworts
(\emph{Chara}) are closest to plants. Like
plants they have chloroplasts with grana,
asymmetrical motile cells with lateral
flagella, a nuclear envelope that breaks down
at mitosis, and a phragmoplast in
cytokinesis; they are oogamous (a non-motile
egg fertilized by a flagellated sperm) and
grow apically. \emph{Coleochaete} forms a
disk of parenchyma one cell thick; its
zygotes stay on the parent thallus, covered
by a layer of parental cells, which in one
species have wall ingrowths like the transfer
cells that feed the embryo of land plants.
This is why Daniel's questionnaire named
\emph{Coleochaete} as a proposed model of the
ancestor of land plants.

Haptophytes. The coccolithophores cover
themselves with scales of calcium carbonate
(coccoliths); they are abundant in marine
plankton and important in the global carbon
budget \cite[Chapter 25]{willey2008prescott}.

\section{Levels of organization and life
cycles}
\label{section-organization}

The levels of organization of algae, as
listed in Daniel's notes and in Lee
\cite[Chapters 2, 4 and 5]{lee2008phycology}:
\begin{itemize}
\item unicellular, motile
(\emph{Chlamydomonas}, \emph{Euglena}) or
not;
\item palmelloid: non-motile cells held
together in mucilage (\emph{Palmella},
\emph{Microcystis});
\item coenobium: a colony with a definite
number of cells arranged in a specific way,
fixed when it forms; cells enlarge but do not
divide afterwards (\emph{Volvox}, the water
net \emph{Hydrodictyon});
\item coenocyte (siphonous thallus): a large
multinucleate cell without cross walls
(\emph{Caulerpa});
\item filament: uniseriate, unbranched or
branched, or polysiphonous (a central
filament surrounded by rows of cells, as in
the red alga \emph{Polysiphonia}), sometimes
covered by a cortex;
\item pseudoparenchyma: interwoven filaments
that look like a tissue;
\item parenchyma: cells dividing in more than
one plane to form a true tissue
(\emph{Coleochaete}, brown algae).
\end{itemize}

Sexual reproduction is \emph{isogamous}
(equal gametes), \emph{anisogamous} (motile
gametes of different size) or \emph{oogamous}
(a large non-motile egg and a small motile
sperm). Life cycles differ in where meiosis
happens \cite[Chapter 12]{evert2013raven}:
\begin{itemize}
\item zygotic meiosis: the zygote is the only
diploid cell and divides by meiosis (the
course's haploid monophasic cycle, as in
\emph{Chlamydomonas});
\item gametic meiosis: the organism is
diploid and the gametes are the only haploid
cells (the diploid monophasic cycle, as in
the brown alga \emph{Fucus});
\item sporic meiosis, alternation of
generations: a diploid sporophyte makes
spores by meiosis, which grow into a haploid
gametophyte (the diphasic cycle); the
generations are isomorphic when they look
alike (\emph{Ulva}) and heteromorphic when
they differ, as in kelps, with a large
sporophyte and a microscopic gametophyte;
\item the triphasic cycle of red algae
described above.
\end{itemize}

\section{Protozoa}
\label{section-protozoa}

Protozoa are unicellular or colonial
eukaryotes without tissues, but with
specialized organelles \cite[Chapter
5]{hickman2002animal} \cite[Chapter
25]{willey2008prescott}.
\begin{itemize}
\item Locomotion: by pseudopodia, flagella or
cilia. Pseudopodia are lobopodia (broad, as
in \emph{Amoeba}), filopodia (thin),
reticulopodia (thin and branching into a net,
as in foraminifera) or axopodia (thin and
stiffened by an axis of microtubules, as in
the sun animalcule \emph{Actinophrys}).
\item Nutrition: autotrophic (with plastids),
holozoic (particles taken by phagocytosis
into food vacuoles) or saprozoic (dissolved
nutrients).
\item Organelles: contractile vacuoles that
expel water in freshwater forms; food
vacuoles; macronucleus and micronucleus in
ciliates; the kinetoplast, a mass of DNA
within the mitochondrion of trypanosomes;
hydrogenosomes, organelles of anaerobic
energy metabolism that oxidize pyruvate and
release $\mathrm{H_2}$ and $\mathrm{CO_2}$,
in the parabasalids instead of mitochondria;
the apical complex of Apicomplexa, with which
they enter host cells.
\item Reproduction: asexual by binary
fission, multiple fission (schizogony) or
budding; sexual by syngamy, by conjugation in
ciliates (exchange of micronuclei), or by
autogamy.
\item Life stages: the active trophozoite and
the resistant cyst, which also spreads
parasites from host to host.
\end{itemize}

Many protozoa live in association with other
organisms. Parabasalids such as
\emph{Trichonympha} live in the gut of
termites and digest the wood they eat; in a
course practical, a toad was dissected to
look for the commensal and parasitic protozoa
of its gut.

\section{Protozoan phyla in the course's
scheme}
\label{section-phyla}

The course classified protozoa after Corliss
(1994) and Aladro Lubel (2006). The class
tables listed these phyla, with
representative genera:

\begin{center}
\begin{tabular}{p{0.4\linewidth}
p{0.5\linewidth}}
Phylum & Genera \\
\hline
Archamoebae & \emph{Pelomyxa},
\emph{Entamoeba} \\
Metamonada & \emph{Giardia} \\
Microspora & \emph{Nosema} \\
Percolozoa & \emph{Naegleria} \\
Parabasala & \emph{Trichomonas},
\emph{Trichonympha} \\
Euglenozoa & \emph{Euglena},
\emph{Leishmania}, \emph{Trypanosoma} \\
Opalozoa & \emph{Opalina} \\
Mycetozoa & \emph{Physarum} \\
Choanozoa & \emph{Codosiga} \\
Dinozoa & \emph{Ceratium} \\
Ciliophora & \emph{Paramecium},
\emph{Vorticella} \\
Apicomplexa & \emph{Monocystis},
\emph{Plasmodium}, \emph{Toxoplasma} \\
Rhizopoda & \emph{Arcella}, \emph{Elphidium}
\\
Heliozoa & \emph{Actinophrys} \\
Radiozoa & \emph{Thalassicolla} \\
Myxozoa & \emph{Myxobolus} \\
Ascetospora & \emph{Haplosporidium} \\
\end{tabular}
\end{center}

Molecular phylogenies have since regrouped
the eukaryotes into a few supergroups; in the
scheme of the International Society of
Protistologists used by Prescott
\cite[Chapter 25]{willey2008prescott}:
\begin{itemize}
\item Excavata: \emph{Giardia} (Fornicata),
the parabasalids, and the Euglenozoa
(euglenoids and trypanosomes);
\item Amoebozoa: naked and testate amoebae,
\emph{Entamoeba}, \emph{Acanthamoeba},
\emph{Pelomyxa} and the slime moulds;
\item Rhizaria: foraminifera, radiolarians
and the haplosporidians;
\item Chromalveolata: the alveolates
(apicomplexans, dinoflagellates, ciliates),
the stramenopiles (diatoms, brown algae,
oomycetes) and the haptophytes;
\item Archaeplastida: glaucophytes, red and
green algae;
\item Opisthokonta: choanoflagellates, with
fungi and animals.
\end{itemize}
The microsporidia, such as \emph{Nosema}, are
now placed among the fungi \cite[Chapter
26]{willey2008prescott}.

\section{Protozoan parasites of humans}
\label{section-parasites}

Daniel's talk began with the kinds of
symbiosis: a parasite depends physiologically
on its host, must find hosts, evade their
immune response and reproduce fast; a
mutualist and its host depend on each other;
a commensal profits without physiological
dependence. Parasites are endo- or
ectoparasites, obligate or facultative (like
\emph{Naegleria}, normally free-living); many
need vectors, have reservoir hosts, or cause
opportunistic infections in immunocompromised
people. The main protozoan parasites, from
the textbooks \cite[Chapter
12]{tortora2013microbiology} \cite[Chapter
39]{willey2008prescott} \cite[Chapter
5]{hickman2002animal}:

\begin{center}
\begin{tabular}{p{0.28\linewidth}
p{0.62\linewidth}}
Parasite & Disease and transmission \\
\hline
\emph{Entamoeba histolytica} & amoebic
dysentery; cysts in faecally contaminated
water. \emph{E. coli}, also common in the
gut, is not a pathogen \\
\emph{Giardia lamblia} & giardiasis, an
enteritis; trophozoite with two nuclei, eight
flagella and a ventral sucking disk; cysts in
contaminated water \\
\emph{Trichomonas vaginalis} & vaginitis and
urethritis; sexually transmitted; no cyst
stage \\
\emph{Naegleria fowleri} & primary amoebic
meningoencephalitis, nearly always fatal
within days; amoeboid, flagellated and cyst
stages; infects swimmers in warm fresh water
through the nasal mucosa \cite[Chapter
22]{tortora2013microbiology}.
\emph{Acanthamoeba} and \emph{Balamuthia}
cause encephalitis and (\emph{Acanthamoeba})
keratitis \\
\emph{Leishmania} & leishmaniasis: visceral
(\emph{L. donovani}, kala-azar), cutaneous
(\emph{L. tropica}, \emph{L. mexicana}),
mucocutaneous (\emph{L. braziliensis});
flagellated promastigotes in female sand
flies (\emph{Phlebotomus}, \emph{Lutzomyia}),
amastigotes inside macrophages; reservoirs in
dogs and rodents \\
\emph{Trypanosoma} & African sleeping
sickness (\emph{T. brucei gambiense} and
\emph{rhodesiense}), by tsetse flies; Chagas
disease (\emph{T. cruzi}) in the Americas, by
triatomine (kissing) bugs whose faeces
contaminate the bite \\
\emph{Plasmodium} & malaria (\emph{P.
falciparum}, \emph{vivax}, \emph{ovale},
\emph{malariae}); sporozoites injected by
female \emph{Anopheles} mosquitoes multiply
in the liver, then merozoites invade red
blood cells in cycles that cause the attacks
of fever; sexual reproduction in the mosquito
\\
\emph{Toxoplasma gondii} & toxoplasmosis,
mild except in AIDS patients and in
congenital infection; the cat is the
definitive host and sheds oocysts in its
faeces; tissue cysts in meat \\
\end{tabular}
\end{center}

The talk also mentioned the ciliate
\emph{Balantidium coli}, microsporidia
(opportunistic in AIDS), \emph{Babesia}
(transmitted by ticks) and
\emph{Cryptosporidium}.

The Mexican context. In Mexico intestinal
protozoa are the commonest parasites. A
survey of Casulco, an irregular settlement in
Magdalena Contreras, Mexico City, with 818
inhabitants in 162 families, found giardiasis
in 29.98\% of those examined and
\emph{Entamoeba histolytica} in 7.29\%, with
the harmless \emph{E. coli} in 14.71\%, and
linked these rates to poor sanitation
\cite{sanchezvega2000frecuencia}. Cutaneous
leishmaniasis by \emph{L. mexicana} occurs in
the Yucat\'an Peninsula \cite[Chapter
39]{willey2008prescott}, and Chagas disease
and trichomoniasis were the other diseases
the talk singled out for Mexico.

\section{The Cantera Oriente field practical}
\label{section-cantera-oriente}

The first practical, on 13 February 2014,
collected algae and protozoa in the Cantera
Oriente, a former basalt quarry on the UNAM
campus that is a buffer zone of the
university's ecological reserve of the
Pedregal de San \'Angel. As described in the
illustrated guide that the course handed out,
its two lakes are fed by canals carrying
water from springs exposed when the rock was
extracted. The team sampled six zones,
recording the water (stagnant or clear),
light, ground and vegetation, and measuring
pH (7.5 to 8.7) and temperature (13 to 17
$^\circ$C). Methods: jars of water with
decaying leaves and floating vegetation,
scraping algae off stones with a toothbrush,
squeezing roots of aquatic plants, and a
plankton net with a 10 $\mu$m mesh, also used
from a boat in the middle of a pond. The
samples were kept refrigerated and examined
as wet mounts under the light microscope.

Organisms identified:
\begin{itemize}
\item diatoms, among them \emph{Cymbella
mexicana} and \emph{Fragilaria};
\item green algae: the desmid
\emph{Closterium}, \emph{Spirogyra} (a
filamentous alga of small stagnant waters,
typical of spring \cite[Chapter
5]{lee2008phycology}), the water net
\emph{Hydrodictyon}, \emph{Cladophora}
(\emph{C. glomerata}) and
\emph{Stichococcus};
\item an ostracod (a small crustacean).
\end{itemize}
Almost all came from one zone with clear
water, solid ground and abundant light, where
samples were taken by scraping and squeezing
roots; the report attributed the empty zones
to samples examined too briefly rather than
to the environment.

\bibliography{my}
\bibliographystyle{amsalpha}

\end{document}
